\runninghead{The Way of the Beatitudes}
\properlane{beatitudes-way-of-the-saints}{The Beatitudes, the Way the Saints Walked}

On the feast of all the saints the Mass reads the eight beatitudes, and at the
Communion it sings the last three of them again. St Augustine reads the
beatitudes as steps that rise from humility, and St John Chrysostom as a
golden chain in which each precept makes way for the next. St Thomas Aquinas
says that their rewards are begun on the way and perfect in the homeland;
Rabanus Maurus, in words he marks as his own, that the clean of heart see God
here by grace and hereafter face to face; St Hilary of Poitiers reads the
promises of those who have been changed by immortality. If the beatitudes are
an ordered way, the Gospel read on this feast describes the road the saints
walked.

\subsection{Why the Beatitudes, and in what tense}
\label{sec:beatitudes-question}

The Gospel does not praise the saints for what they now enjoy. It calls
blessed the poor in spirit, the meek, those who mourn, those who hunger and
thirst, those who are persecuted and reviled: people who lack, grieve and
suffer. Its promises raise a second difficulty, of time. The first and the
eighth beatitude say \latin{ipsorum est regnum caelorum}, ``theirs is the
kingdom''; the six between say that the blessed shall possess the land, be
comforted, have their fill, obtain mercy, see God, be called the children of
God; and the last verse says \latin{merces vestra copiosa est in caelis},
``your reward is very great in heaven.'' If every promise waits for the life to
come, the saints on earth had none of what was promised; if every promise is
given now, the mourners are comforted already, which their mourning does not
show. A third question is the Communion: why the Mass, having read eight
beatitudes, sings three of them again.

\subsection{St Augustine: seven steps and an eighth}
\label{sec:beatitudes-augustine}

St Augustine reads the eight sayings as stages in order, rising from humility
(\work{De serm.\ Dom.}\ I.3.10). The poor in spirit are the humble and those
who fear God; the soul that has submitted to God becomes meek, learns in what
evils it is held and mourns, hungers for justice, shows mercy, is purified in
heart, and at the seventh step reaches wisdom, the peace of the whole man.
``The eighth, as it were, returns to the starting-point, because it shows and
commends what is complete and perfect: therefore in the first and in the
eighth the kingdom of heaven is named.'' He then sets the seven steps beside
the seven operations of the Holy Ghost in Isaias, counted upward from the
fear of the Lord (Isa 11:2--3): fear answers to the poor in spirit, piety to
the meek, knowledge to the mourners, fortitude to the hungry, counsel to the
merciful, understanding to the clean of heart, and wisdom to the peacemakers
(I.4.11). Augustine puts the fear of the Lord at the first step, beside the
poor in spirit; if the Gradual's \latin{Timete Dominum} is that fear, the
Gradual sings the first step.

The rewards, he says, are one reward, the kingdom of heaven, variously named
according to the stages, and the eighth ``goes back to the starting-point, and
makes manifest the perfect man,'' as circumcision was on the eighth day and the
Lord rose on the day after the Sabbath (I.4.12). Of the time of the rewards
he writes: ``And those promises can indeed be fulfilled in this life, as we
believe them to have been fulfilled in the case of the apostles. For that
all-embracing change into the angelic form, which is promised after this
life, cannot be explained in any words.'' Preaching on the Alleluia's verse
and the one after it, ``learn of me, because I am meek, and humble of heart,''
he tells his hearers where such a way begins: ``Thou wishest to be great, begin
from the least,'' and ``first think of the foundation of humility''
(\work{Sermon} XIX [LXIX], n.~2).

\subsection{St John Chrysostom: a golden chain, and comfort here and there}
\label{sec:beatitudes-chrysostom}

St John Chrysostom also reads the beatitudes in order. The Lord, he says,
``by the former precept making way for the following one, He hath woven a
sort of golden chain for us'': the humble man will mourn for his sins; the
mourner will be meek, righteous and merciful; the merciful will be pure in
heart and a peacemaker; and the man who has attained all these is armed
against dangers and is not troubled when evil is spoken of him
(\work{Hom.\ in Matt.}\ 15.9). He asks where the mourners shall be
comforted, and answers, ``Both here and there'' (15.4). Of the meek who
inherit the earth he says that the Lord speaks ``not as limiting the rewards
to things present, but as joining with these the other sort of gifts also''
(15.5).

On the Alleluia's verse he makes the link to the Gospel himself. The meekness
and lowliness of which the Lord speaks in the next verse, Mt 11:29, are ``the
mother of all strictness of life. Wherefore also, when beginning those divine
laws, with this He began,'' that is, with the Beatitudes; and of the rest
promised there to the meek and lowly, ``even before the things to come, He
gives thee here thy recompense'' (\work{Hom.\ in Matt.}\ 38.3). Chrysostom
and Augustine both begin the way in humility, and both read the Alleluia's
call with the verse that follows it.

\subsection{St Thomas and Rabanus Maurus: begun in the way, perfect in the homeland}
\label{sec:beatitudes-thomas}

St Thomas Aquinas gives the distinction in two words. The rewards are had
\latin{in via}, on the way, as a beginning (\latin{secundum inchoationem}),
and \latin{in patria}, in the homeland, perfectly (\work{Super Matt.}\ 5,
Venice printing, p.~67). Of the peacemakers he says that their peace
\latin{hic inchoatur}, ``is begun here'' (p.~72), and the eight rewards are
\latin{unum praemium}, one reward under different names. On the eighth he
reports Chrysostom: \latin{Non dicit a paganis}, the Lord does not say ``by the
pagans,'' since persecution for justice is suffered not only from pagans. The
Leonine text of the same lectures adds that the kingdom is promised to the
poor in hope and to the martyrs in fact.

Rabanus Maurus compiles his commentary on Matthew from earlier writers and
marks many of his sources in the running text: ``(Hieron.)'' for Jerome,
``(Aug.)'' for Augustine, ``(MAURUS.)'' for himself. From Augustine he takes the
poor in spirit as those who fear God, \latin{Initium autem sapientiae timor
Domini}, ``the beginning of wisdom is the fear of the Lord,'' and, at 5:9--10,
Augustine's sentence that the promises can be fulfilled in this life and that
the eighth returns to the start; at the Lord's Prayer he repeats Augustine's
matching of the seven petitions to the beatitudes. Under his own name he says
of the clean of heart: \latin{Dominum enim mundi corde vident hic in aenigmate
per gratiam, in futuro autem facie ad faciem}, ``the clean of heart see the
Lord here in a dark manner by grace, and in the future face to face'' (on 5:8,
with 1~Cor 13:12). St Bernard, preaching on the feast, says the same of the
vision: \latin{Videbunt quidem in fine facie ad faciem; videbunt etiam nunc,
sed per speculum in aenigmate}, ``they will see at the end face to face; they
see even now, but through a glass in a dark manner'' (sermo~1, n.~13), and he
calls the lesser kingdom given now a pledge, \latin{velut in arrham}, of the
greater (n.~9).

\subsection{St Hilary: the precepts of the heavenly life}
\label{sec:beatitudes-hilary}

St Hilary of Poitiers reads the promises chiefly of the state the saints now
possess. The Lord on the mountain, \latin{in paternae scilicet majestatis
positus celsitudine}, ``set, that is, in the height of the Father's majesty,''
lays down \latin{caelestis vitae praecepta}, the precepts of the heavenly life,
for he would not have handed on the laws of eternity had he not stood in
eternity (\work{In Matt.}\ IV, n.~1). The earth the meek inherit is the body
the Lord assumed, and \latin{claritate clarificati ejus corporis vestiemur},
``we shall be clothed with the glory of his glorified body'' (n.~3).
Consolation and fullness are \latin{in caelo}, in heaven (nn.~4--5). The clean
of heart see God because \latin{non enim nisi spiritu perfecti, et
immortalitate immutati \dots\ hoc quod in Deo est immortale cernemus}, ``only
when we are perfected in spirit and changed by immortality shall we behold what
in God is immortal'' (n.~7). The persecuted are \latin{gloriosi promissorum Dei
martyres}, the glorious martyrs of God's promises (n.~9).

\subsection{Joy, peace, and three beatitudes sung again}
\label{sec:beatitudes-elements}

The Introit's verse bids the just rejoice, and St Basil, on that verse,
sets their joy beside ours: the joy God sometimes, rarely, lets fall on us is
a small taste that reminds us of what we lack, while \latin{justo perpetua est
divina ac caelestis laetitia, quod semel in ipso inhabitat Spiritus}, ``the
just man has a divine and heavenly joy that does not end, because the Spirit
dwells in him once for all'' (\work{Hom.\ in Ps.}\ 32, n.~1).

The Offertory says of the just, \latin{illi autem sunt in pace}, ``but they are
in peace.'' St Bonaventure, on these words, reports the \work{Glossa}:
\latin{Perpetuae quietis, nunc in spe, tandem in re}, ``of perpetual rest, now
in hope, at last in fact.'' St Alphonsus Liguori, on the same verses, says that
the saints at death ``begin to feel that peace which they shall afterwards
fully enjoy in Heaven.'' Fromage, on the Offertory, joins it to the Gospel just
read: the Beatitudes raise man above torments and above death itself, which
\latin{n'ébranle pas la paix des justes, mais la consomme}, does not shake the
peace of the just but consummates it, and that, he says, is what the Offertory
sings.

The Communion sings the clean of heart, the peacemakers and the persecuted,
word for word as the Gospel has just given them. Fromage gives the
reason for the three: \latin{Echo de la lecture évangélique, l'Antienne de
Communion, ne pouvant énumérer à nouveau la série entière des Béatitudes,
rappelle les trois dernières, et, ce faisant, les rapproche toutes avec raison
du Sacrement divin où elles s'alimentent.} The antiphon echoes the Gospel;
since it cannot list the whole series again, it recalls the last three, and in
doing so brings them all, rightly, to the divine Sacrament in which they are
nourished. The Fathers above comment on the same three verses as part of the
Gospel. In Augustine's order they are the sixth and seventh steps and the
eighth that returns to the first. St Bernard's first sermon for the feast expounds the Gospel's
beatitudes as reconciliations: the first three of the soul with itself, the
next two with its neighbour, the sixth with God; in the seventh it reconciles
others, and the eighth is the martyrs' prerogative (sermo~1, nn.~14--15).

The Lesson's multitude wears white robes and has palms in its hands. Bede
reads them first as those who overcome the last persecution (first
interpretation above); if they are, they are the persecuted of the eighth
beatitude, whose reward is the kingdom of heaven. The feast's three orations
venerate the saints' merits and ask mercy through their many intercessors,
offer the gifts for the honour of all the just, and ask that the faithful ever
rejoice in venerating the saints. The Sunday's Collect asks release from sins contracted
\latin{pro nostra fragilitate}, through our frailty; its Secret, that God
complete what he has given to the undeserving; its Postcommunion, protection
from human dangers. Either Preface ends with the praise of the angels.

\subsection{Present and future}
\label{sec:beatitudes-difficulty}

The tenses of the Gospel are the hardest point. The first and the eighth
beatitude are in the present, \latin{est}, and the six between in the future.
The Fathers do not divide the rewards by those tenses. Augustine calls the
eight rewards one reward, the kingdom of heaven, variously named, and Aquinas
calls them \latin{unum praemium}; what they say of time, they say of all the
rewards together.

Augustine, Chrysostom, Aquinas and Rabanus place the rewards in this life,
as fulfilment, comfort, beginning or vision by grace. Augustine says that the
promises can be fulfilled here, as they were in the apostles; Chrysostom, that
the mourners are comforted here as well as there, and that the promise of the
earth does not limit the rewards to the present; Aquinas, that peace is begun
here; Rabanus, that the clean of heart see the Lord here by grace. Bernard's
pledge and mirror say the same. Hilary places the rewards in heaven, for those
changed by immortality. The difference is real, and it changes what the Gospel
describes: on Hilary's account chiefly the state the saints now possess, on the
others' the life they already lived. Augustine's own next sentence keeps a
change after this life that no words can explain, so between him and Hilary the
difference is one of emphasis.

Bernard gives a reason why the first and the eighth have the same promise. The
kingdom is promised to the poor and to the martyrs alike, \latin{quia vere
martyrii genus paupertas voluntaria est}, because voluntary poverty is truly a
kind of martyrdom: the kingdom \latin{paupertate quidem emitur, sed in
passione pro Christo absque omni dilatione percipitur}, ``is bought by poverty,
but received without delay in suffering for Christ'' (sermo~1, n.~15).

\begin{fourSenses}
\item[Literal.] Jesus, seated on a mountain, teaches his disciples that the
poor in spirit, the meek, the mourners, those who hunger for justice, the
merciful, the clean of heart, the peacemakers and the persecuted are blessed,
and tells the disciples that when they are reviled and persecuted for his sake
their reward is great in heaven.
\item[Allegorical.] Christ teaches from the height of the Father's majesty the
precepts of the heavenly life, and the earth of the meek is his glorified
body, with whose glory the saints will be clothed (Hilary).
\item[Moral.] The beatitudes are steps rising from humility and the fear of
the Lord, answering to the gifts of the Holy Ghost (Augustine), and a golden
chain in which each precept makes way for the next, from the humble to the
persecuted (Chrysostom); they reconcile the soul with itself, with its
neighbour and with God (Bernard).
\item[Anagogical.] The one reward, begun on the way, is given perfectly in the
homeland (Aquinas), to those perfected in spirit and changed by immortality,
who see what in God is immortal (Hilary); the change into the angelic form is
promised after this life (Augustine).
\end{fourSenses}
